% TL;DR: Projektkarte als Poster A1 quer (84 × 59,4 cm): U-Bahn-Plan, Ziel und Pipeline, Zahlen, Wendepunkte, Sackgassen, Lessons Learned, Ausblick; braucht build/karte_voll.pdf aus build.sh.
\documentclass{article}
\usepackage[paperwidth=841mm, paperheight=594mm, margin=0mm]{geometry}
\usepackage{libertinus-otf}
\usepackage{xcolor}
\usepackage{graphicx}
\usepackage{tikz}
\usepackage[most]{tcolorbox}
% TL;DR: Gemeinsame Farben, TikZ-Stile und Stations-Makros der Projektkarte (U-Bahn-Plan); wird von karte.tex und den Varianten eingebunden.
%
% Lesart der Karte
%   Linie       = Arbeitsstrang (Farbe)
%   Station     = Meilenstein; die Füllung zeigt den Ausgang (\Mk)
%   Stumpfgleis = aufgegebener Weg (endet mit Prellbock)
%   Fehlannahme = gestrichelte Nebenspur (Stil irrweg): so lange galt eine falsche Annahme

\usetikzlibrary{positioning, calc, patterns, shapes.geometric, fit, backgrounds, arrows.meta}

% ── Linienfarben (Arbeitsstränge) ─────────────────────────────────────
\definecolor{lInfra}{HTML}{6B7B8C}   % Infrastruktur / Rechenzugang
\definecolor{lTrain}{HTML}{1F4E9C}   % Training N1.6 (NavyBlue wie die Ausarbeitung)
\definecolor{lNsieben}{HTML}{8E44AD} % GR00T N1.7
\definecolor{lCotr}{HTML}{D1477A}    % Co-Training / Daten
\definecolor{lSim}{HTML}{0E8C7F}     % Simulation / Szene
\definecolor{lDiag}{HTML}{E07B00}    % Evaluation / Diagnose
\definecolor{lRL}{HTML}{8D8D8D}      % Reinforcement Learning
\definecolor{lDoku}{HTML}{3C3C3C}    % Doku / Abgabe

% ── Statusfarben (Stationen) ──────────────────────────────────────────
\definecolor{sOk}{HTML}{2E8B3E}      % Erfolg
\definecolor{sFail}{HTML}{C62828}    % Fehlschlag
\definecolor{sPart}{HTML}{F2B705}    % Teilerfolg
\definecolor{sOld}{HTML}{9E9E9E}     % verworfen / überholt
\definecolor{phaseA}{HTML}{F4F6F9}
\definecolor{phaseB}{HTML}{E9EDF3}
\definecolor{hubbg}{HTML}{FFF8E1}

% ── Maße (die Varianten skalieren die fertige Karte als Ganzes) ───────
\newcommand\linew{2.6pt}
\newcommand\fsLbl{\fontsize{6.4}{7.3}\selectfont\sffamily}
\newcommand\fsPhase{\fontsize{7.2}{8.4}\selectfont\sffamily}

\tikzset{
  line/.style   = {line width=\linew, line cap=round, line join=round, rounded corners=4mm},
  stub/.style   = {line width=\linew, line cap=butt, dash pattern=on 3pt off 2pt},
  irrweg/.style = {line width=1.1pt, dash pattern=on 2.4pt off 1.6pt, line cap=round,
                   rounded corners=2.5mm},
  lbl/.style    = {font=\fsLbl, align=center, inner sep=0.6pt, text depth=0pt},
  badge/.style  = {rounded corners=1.6mm, font=\fontsize{6.8}{7.6}\selectfont\sffamily\bfseries,
                   text=white, inner xsep=1.6mm, inner ysep=1.1mm, minimum width=19mm,
                   align=center},
}

% \Mk{status}{x,y}: Stationsmarke. status ∈ ok | fail | part | old | open
\newcommand\Mk[2]{%
  \begin{scope}[shift={(#2)}]
    \fill[white] (0,0) circle (1.95mm);
    \csname mk@#1\endcsname
  \end{scope}}
\makeatletter
\newcommand\mk@ok{\fill[sOk] (0,0) circle (1.45mm);
  \draw[white, line width=0.85pt, line cap=round, line join=round]
    (-0.72mm,0.02mm) -- (-0.18mm,-0.58mm) -- (0.78mm,0.62mm);}
\newcommand\mk@fail{\fill[sFail] (0,0) circle (1.45mm);
  \draw[white, line width=0.85pt, line cap=round]
    (-0.6mm,-0.6mm) -- (0.6mm,0.6mm) (-0.6mm,0.6mm) -- (0.6mm,-0.6mm);}
\newcommand\mk@part{\fill[white] (0,0) circle (1.45mm);
  \fill[sPart] (0,-1.45mm) arc (-90:90:1.45mm) -- cycle;
  \draw[sPart!80!black, line width=0.6pt] (0,0) circle (1.2mm);}
\newcommand\mk@old{\fill[sOld] (0,0) circle (1.45mm);
  \draw[white, line width=0.8pt] (-0.75mm,0) -- (0.75mm,0);}
\newcommand\mk@open{\fill[white] (0,0) circle (1.45mm);
  \draw[black!70, line width=0.6pt, dash pattern=on 0.9pt off 0.7pt] (0,0) circle (1.2mm);}
\makeatother

% \Wende{x,y}: Wendepunkt (schwarzer Ring um die Station)
\newcommand\Wende[1]{\draw[black, line width=0.9pt, fill=white] (#1) circle (2.55mm);}

% \St[*]{status}{x}{y}{Platzierung}{Titel}{Detail}
%   *            = Wendepunkt
%   Platzierung  = above | below | left | right | above left | ... (Abstand fest)
% Abstand der Beschriftung zur Marke; Wendepunkte brauchen mehr wegen des Rings
\newcommand\lblsep{2.35mm}
\newcommand\lblsepW{3.3mm}
\NewDocumentCommand\St{s m m m m m m}{%
  \IfBooleanTF{#1}{\Wende{#3,#4}\def\sep{\lblsepW}}{\def\sep{\lblsep}}%
  \Mk{#2}{#3,#4}%
  \node[lbl, #5=\sep] at (#3,#4) {{\bfseries #6}\ifx\relax#7\relax\else\\#7\fi};}

% \Prellbock{x,y}{Winkel}: Ende eines Stumpfgleises (Querbalken)
\newcommand\Prellbock[3][black]{%
  \draw[#1, line width=1.6pt, line cap=butt] ($(#2)+(#3+90:1.7mm)$) -- ($(#2)+(#3-90:1.7mm)$);}

% \Legende: Zeichenerklärung als eigenes Bild (Varianten skalieren mit \scalebox)
\newcommand\Legende{%
\begin{tikzpicture}[x=1cm, y=1cm, every node/.style={font=\fsLbl, anchor=west, inner sep=1pt}]
  \Mk{ok}{0,0}\node at (0.25,0) {Erfolg};
  \Mk{part}{1.45,0}\node at (1.7,0) {Teilerfolg};
  \Mk{fail}{3.25,0}\node at (3.5,0) {Fehlschlag};
  \Mk{old}{5.1,0}\node at (5.35,0) {verworfen / überholt};
  \Wende{7.95,0}\Mk{ok}{7.95,0}\node at (8.25,0) {Wendepunkt};
  \draw[stub, black!70] (10.1,0) -- (10.9,0); \Prellbock[black!70]{10.9,0}{0}
  \node at (11.05,0) {Stumpfgleis = aufgegebener Weg};
  \draw[irrweg, black!70] (15.3,0) -- (16.1,0);
  \node at (16.2,0) {Fehlannahme = Dauer bis zur Widerlegung};
  \node at (20.55,0) {T$n$ = Trainingslauf $n$ · L$n$ = Diagnoselauf $n$ · Zeitachse nicht maßstäblich};
\end{tikzpicture}}

% \BlockerKette: Ursachenkette als Zeile (jede gelöste Ebene legte die nächste frei)
\newcommand\BlockerKette{%
\begin{tikzpicture}[x=1cm, y=1cm,
    bk/.style={draw=black!45, line width=0.5pt, fill=white, rounded corners=1.2mm,
      font=\fsLbl, align=center, text width=31mm, minimum height=9.5mm, inner sep=1.2pt},
    pf/.style={-{Latex[length=1.6mm]}, line width=0.9pt, black!55}]
  \node[font=\fsLbl\bfseries, anchor=east, align=right] at (-0.25,0)
    {Blocker-Kette:\\[-1pt]{\mdseries jede gelöste Ebene}\\[-1pt]{\mdseries legte die nächste frei}};
  \foreach \i/\s/\t/\d in {
      0/ok/{Kameras verdreht}/{behoben, L10–18},
      1/ok/{Geometrie: Finger, Becken, Basis}/{behoben, L38–51},
      2/part/{Würfel 3 cm versetzt}/{korrigiert, Ursache offen, L62},
      3/part/{Bild-Lücke Sim ↔ Echt}/{×13 → halbiert, L66/L83},
      4/part/{Griff}/{2/60 frei gehoben, L97},
      5/open/{Stapeln + Robotertest}/{0/60, offen}} {
    \node[bk] (b\i) at ({\i*3.75+1.75},0) {{\bfseries\t}\\\d};
    \Mk{\s}{{\i*3.75+0.12},0.47}
  }
  \foreach \i [evaluate=\i as \j using int(\i+1)] in {0,...,4} \draw[pf] (b\i.east) -- (b\j.west);
\end{tikzpicture}}

\pagestyle{empty}
\setlength{\parindent}{0pt}

% ── Poster-Typografie ─────────────────────────────────────────────────
\newcommand\pBody{\fontsize{19}{24}\selectfont\sffamily}
\newcommand\pSmall{\fontsize{16.5}{21}\selectfont\sffamily}
\newcommand\pBig{\fontsize{30}{34}\selectfont\sffamily\bfseries}

\newtcolorbox{panel}[4][]{enhanced, width=#2, title={#4}, sharp corners=south, arc=3mm,
  colback=white, colframe=#3, boxrule=1.2mm, toprule=0mm, titlerule=0mm,
  colbacktitle=#3, coltitle=white,
  fonttitle=\fontsize{24}{28}\selectfont\sffamily\bfseries,
  left=5mm, right=5mm, top=3.5mm, bottom=3mm, toptitle=2mm, bottomtitle=2mm, #1}

\newcommand\Pkt[2][lTrain]{\textcolor{#1}{\rule[0.22em]{0.42em}{0.42em}}\hspace{0.45em}#2}

\newsavebox\legbox \newsavebox\bkbox
\begin{document}
\sbox\legbox{\Legende}\sbox\bkbox{\BlockerKette}
\noindent
\begin{tikzpicture}[remember picture, overlay, x=1mm, y=1mm]
  \coordinate (O) at (current page.north west);

  % ── Kopf ──────────────────────────────────────────────────────────
  \fill[lTrain] (O) rectangle ($(O)+(841,-80)$);
  \node[anchor=west, fill=white, rounded corners=3mm, inner sep=4mm] at ($(O)+(15,-40)$)
    {\includegraphics[height=46mm]{../logo_hochschule.jpg}};
  \node[anchor=north west, text=white, align=left, inner sep=0] at ($(O)+(128,-13)$) {%
    {\fontsize{50}{56}\selectfont\sffamily\bfseries Projektkarte: GR00T-Fine-Tuning für Unitree G1 + DEX3}\\[4mm]
    {\fontsize{26}{31}\selectfont\sffamily Entwicklungsweg Mai bis November 2026 – Arbeitsstränge,
      Erfolge, Fehlschläge und Irrwege bis zum heutigen Stand}\\[5mm]
    {\fontsize{21}{25}\selectfont\sffamily Maximilian Berger · Matthias Späth · Luca Mücke
      \quad|\quad Betreuer: Prof. Dr. Matthias Lorenzen
      \quad|\quad Hochschule für angewandte Wissenschaften Kempten}};
  \node[anchor=north east, text=white!80!lTrain, align=right, inner sep=0,
    font=\fontsize{17}{21}\selectfont\sffamily] at ($(O)+(826,-13)$)
    {Projektarbeit 2026\\Stand 01.11.2026};

  % ── Karte ─────────────────────────────────────────────────────────
  \node[anchor=north west, inner sep=0] at ($(O)+(12,-90)$)
    {\includegraphics[width=648mm]{build/karte_voll.pdf}};

  % ── Blocker-Kette + Legende unter der Karte ───────────────────────
  \node[anchor=north, inner sep=0] at ($(O)+(336,-472)$) {\scalebox{2.08}{\usebox\bkbox}};
  \node[anchor=south, inner sep=0] at ($(O)+(420.5,-588)$) {\scalebox{2.3}{\usebox\legbox}};

  % ── Rechte Spalte ─────────────────────────────────────────────────
  \node[anchor=north west, inner sep=0] at ($(O)+(672,-90)$) {%
  \begin{minipage}{157mm}
  \begin{panel}{157mm}{lTrain}{Worum geht es?}\pBody\raggedright
    NVIDIA GR00T, ein Vision-Language-Action-Modell, lernt aus echten Demonstrationen,
    mit dem Unitree G1 und der DEX3-Hand Würfel zu stapeln. Geprüft wird in einer
    Isaac-Lab-Simulation im geschlossenen Regelkreis.\par\smallskip
    \centering
    \begin{tikzpicture}[x=1mm, y=1mm, every node/.style={font=\pSmall, align=center,
        rounded corners=2mm, draw=black!40, line width=0.5mm, fill=phaseA, text width=128mm,
        inner sep=2.2mm}]
      \node (a) at (0,0)   {Echte Demos\\{\color{black!60}unitreerobotics/G1\_Dex3\_BlockStacking}};
      \node (b) at (0,-19.5) {Fine-Tuning im Docker-Image\\{\color{black!60}KISSKI, 4× A100, 44\,000 Schritte}};
      \node (c) at (0,-39) {Checkpoint\\{\color{black!60}Hugging Face, Open-Loop-Val-MSE}};
      \node (d) at (0,-58.5) {Closed Loop in Isaac Lab\\{\color{black!60}Institutsserver, RT-Core-GPU}};
      \node (e) at (0,-78) {erreicht → berührt → gehoben → gestapelt};
      \foreach \u/\v in {a/b,b/c,c/d,d/e}
        \draw[-{Latex[length=4mm]}, line width=0.8mm, black!55, draw=black!55] (\u) -- (\v);
    \end{tikzpicture}
  \end{panel}
  \vspace{4mm}
  \begin{panel}{157mm}{lDiag}{Projekt in Zahlen}\pBody
    \renewcommand{\arraystretch}{1.0}
    \begin{tabular}{@{}r@{\hspace{4mm}}l@{}}
      {\pBig 6} & Monate, 08.05.–01.11.2026\\
      {\pBig 301} & Commits\\
      {\pBig 97} & protokollierte Läufe\\
      {\pBig 8} & Trainingsläufe (4 + 4 Co-Training)\\
      {\pBig 2} & GR00T-Generationen\\
      {\pBig 3} & Rechenorte\\
      {\pBig 2/60} & erste freie Griffe in der Sim\\
    \end{tabular}
  \end{panel}
  \vspace{4mm}
  \begin{panel}{157mm}{lSim}{Lessons Learned}\pSmall\raggedright
    \Pkt[lSim]{\textbf{Referenz zuerst:} erst muss die abgespielte Demo in der Sim greifen.}\par\smallskip
    \Pkt[lSim]{\textbf{Wirkung messen, nicht Ähnlichkeit:} SigLIP 0,22 sah gut aus, Fehler ×13.}\par\smallskip
    \Pkt[lSim]{\textbf{Szene unabhängig prüfen:} gegen Echtbilder, nicht mit demselben Code.}\par\smallskip
    \Pkt[lSim]{\textbf{Genug Episoden, feste Seeds:} 6×300 lieferte ein Scheinoptimum.}\par\smallskip
    \Pkt[lSim]{\textbf{Regeln vorab als Zahl:} das vorregistrierte Gate verhinderte Umdeuten.}\par\smallskip
    \Pkt[lSim]{\textbf{Laufenden Prozess prüfen:} ein Schalter in der Konfiguration ist kein Beleg.}\par\smallskip
    \Pkt[lSim]{\textbf{Hardware je Teilaufgabe klären:} Rendern braucht RT-Cores.}\par\smallskip
    \Pkt[lSim]{\textbf{Laborbuch führen:} Laufnummern machten Korrekturen nachverfolgbar.}
  \end{panel}
  \end{minipage}};

  % ── Untere Reihe ──────────────────────────────────────────────────
  \node[anchor=north west, inner sep=0] at ($(O)+(12,-506)$) {%
  \begin{panel}[height=68mm]{250mm}{black!75}{Wendepunkte}\pSmall
    \renewcommand{\arraystretch}{1.0}%
    \begin{tabular}{@{}l@{\hspace{3mm}}l@{}}
      13.08. & span-Gate: echt 100\,\%, Sim 19\,\% Fingerspanne (L32)\\
      24.08. & Geometrie: die Hand konnte nie schließen (L38–51)\\
      12.10. & kein Overfitting, @44\,000 ist am besten\\
      14.10. & Demo hebt kaum: Würfel 3 cm versetzt (L57/62)\\
      14.10. & Sim-Bilder so schlecht wie schwarze (L65/66)\\
      29.10. & Gate bestanden, erste Griffe 2/60 (L96/97)\\
    \end{tabular}
  \end{panel}};
  \node[anchor=north west, inner sep=0] at ($(O)+(271,-506)$) {%
  \begin{panel}[height=68mm]{290mm}{sFail}{Sackgassen \& Irrwege}\pSmall
    \begin{minipage}[t]{137mm}
      \Pkt[sFail]{Lauf 2: Encoder trainiert → Kollaps}\\
      \Pkt[sFail]{Isaac Sim 5.1 stürzt ab → 6.0}\\
      \Pkt[sFail]{Sim auf KISSKI: keine RT-Cores}\\
      \Pkt[sFail]{cotrain\_synth\_v1: nur 1 Datensatz}\\
      \Pkt[sOld]{Cosmos: nie Ende-zu-Ende}
    \end{minipage}\hfill
    \begin{minipage}[t]{137mm}
      \Pkt[sOld]{Finger-Vorzeichen: Hand schloss nie}\\
      \Pkt[sOld]{SigLIP-Abstand als Gap-Maß}\\
      \Pkt[sOld]{U-Kurve @30\,000: Artefakt}\\
      \Pkt[sOld]{Messpunkt 3× falsch (L24–28)}\\
      \Pkt[lRL]{RL: ohne Griff kein Lernsignal}
    \end{minipage}
  \end{panel}};
  \node[anchor=north west, inner sep=0] at ($(O)+(570,-506)$) {%
  \begin{panel}[height=68mm]{259mm}{lCotr}{Ausblick}\pSmall
    \Pkt[lCotr]{Test am echten G1: co-trainiertes cube4 gegen vision-v2}\\
    \Pkt[lCotr]{mehr Render-Episoden statt höherer Mischung}\\
    \Pkt[lCotr]{RL auf dem co-trainierten Modell (erstes Griffsignal)}\\
    \Pkt[lCotr]{Gate-Schwelle relativ statt 80\,\%, Kontaktkriterium für Griffe}\\
    \Pkt[lCotr]{Ursache des 3-cm-Versatzes und der Hänger klären}
  \end{panel}};
\end{tikzpicture}
\end{document}
